\documentclass{tufte-handout}

%\geometry{showframe}% for debugging purposes -- displays the margins

\usepackage[spanish, es-tabla]{babel}
\usepackage{amsmath}

% Set up the images/graphics package
\usepackage{graphicx}
\setkeys{Gin}{width=\linewidth,totalheight=\textheight,keepaspectratio}
\graphicspath{{graphics/}}

\title[Química Física II: Ejercicios Tema 6]{
Química Física II: Ejercicios Tema 6}
%\author{David De Sancho}
\date{}  % if the \date{} command is left out, the current date will be used

% The following package makes prettier tables.  We're all about the bling!
\usepackage{booktabs}

% The units package provides nice, non-stacked fractions and better spacing
% for units.
\usepackage{units}

% The fancyvrb package lets us customize the formatting of verbatim
% environments.  We use a slightly smaller font.
\usepackage{fancyvrb}
\fvset{fontsize=\normalsize}

% Small sections of multiple columns
\usepackage{multicol}

% Provides paragraphs of dummy text
\usepackage{lipsum}

% These commands are used to pretty-print LaTeX commands
\newcommand{\doccmd}[1]{\texttt{\textbackslash#1}}% command name -- adds backslash automatically
\newcommand{\docopt}[1]{\ensuremath{\langle}\textrm{\textit{#1}}\ensuremath{\rangle}}% optional command argument
\newcommand{\docarg}[1]{\textrm{\textit{#1}}}% (required) command argument
\newenvironment{docspec}{\begin{quote}\noindent}{\end{quote}}% command specification environment
\newcommand{\docenv}[1]{\textsf{#1}}% environment name
\newcommand{\docpkg}[1]{\texttt{#1}}% package name
\newcommand{\doccls}[1]{\texttt{#1}}% document class name
\newcommand{\docclsopt}[1]{\texttt{#1}}% document class option name

\begin{document}

\maketitle
\large 
\begin{enumerate}
    \item Una partícula en una caja de longitud $a$ se encuentra sometida a un
    potencial $V_0x/a$ entre $0$ y $a$. Para aproximar la función de onda
    del estado fundamental y su energía podemos usar el método variacional 
    lineal usando las funciones propias de la partícula en una caja para $n=1$
    y $n=2$. Escribe el determinante secular que permitirá encontrar la mejor 
    aproximación de la energía.
    \item Una partícula se encuentra sometida al siguiente potencial
    \begin{equation*}
\hat{V}(x)=
\begin{cases}
  \infty & \text{ si } L<x<0\\
  0  & \text{ si } 0\leq x\leq \frac{L}{4}\text{  o  }\frac{3L}{4}\leq x\leq L
\\
  V_0, & \text{si } \frac{L}{4}<x<\frac{3L}{4}
\end{cases}
\end{equation*}
    donde $V_0$ es una constante pequeña. Usa la teoría de perturbaciones
    para calcular las correcciones de primer orden de la energía, utilizando
    como hamiltoniano de orden cero el correspondiente a la partícula en 
    una caja.
    
    \item Calcula la corrección de primer orden a la energía del 
    estado fundamental de un oscilador armónico cuya energía potencial 
    es 
    \begin{equation*}
        V(x)=\frac{1}{2}kx^2 + \frac{1}{6}\gamma x^3  + \frac{1}{24}bx^4
    \end{equation*}
\end{enumerate}

\newpage
\subsection{\textbf{Soluciones:}}
\begin{enumerate}
\item 
\begin{equation}\begin{vmatrix}
  \frac{h}{8ma^2} + \frac{V_0}{2} - E & -\frac{16V_0}{9\pi^2} \\
 -\frac{16V_0}{9\pi^2} &  \frac{4h^2}{8ma^2} + \frac{V_0}{2} - E 
\end{vmatrix} = 0
\end{equation}
\item $E_n^{(1)}=\frac{V_0}{2}  - \frac{V_0}{2n\pi}(\sin(3n\pi/2) - \sin(n\pi/2))$
\item $E_{v=0}^{(1)}= \frac{b\hbar^2}{32km}$; $E=E_{v=0}^{(0)}+E_{v=0}^{(1)}= 1/2\hbar\omega +  \frac{b\hbar^2}{32km}$.
\end{enumerate}

\end{document}